% from ega2-1-fr.tex
\oldpage[II]{26}
pour tout $n>0$~: en effet, si
$\mathfrak{p}\in\operatorname{Proj}(S)$ contient les éléments homogènes de
$\mathfrak{J}$ de degré $\geq n$, comme par hypothèse il existe un élément
homogène $f\in S_d$ non contenu dans $\mathfrak{p}$, pour tout $m\geq 0$ et
tout $x\in S_m\cap\mathfrak{J}$, on a
$f^r x\in\mathfrak{J}\cap S_{m+rd}$ sauf pour un nombre fini de valeurs de
$r$, donc $f^r x\in\mathfrak{p}\cap S_{m+rd}$, ce qui entraîne
$x\in\mathfrak{p}\cap S_m$
(\hyperref[II.2.1.9-fr]{2.1.9}).

Enfin, on a, pour tout idéal gradué $\mathfrak{J}$ de $S$
\[
  \label{II.2.3.2.5-fr}
  V_+(\mathfrak{J})=V_+(r_+(\mathfrak{J})).
  \tag{2.3.2.5}
\]
\end{env}

\begin{env}[2.3.3]
\label{II.2.3.3-fr}
Par définition, les $V_+(E)$ sont les parties fermées de
$X=\operatorname{Proj}(S)$ pour la topologie induite par la topologie
spectrale de $\operatorname{Spec}(S)$, et qu'on appellera encore
\emph{topologie spectrale} sur $X$. On pose, pour tout $f\in S$
\[
  \label{II.2.3.3.1-fr}
  D_+(f)=D(f)\cap\operatorname{Proj}(S)
  =\operatorname{Proj}(S)\setminus V_+(f)
  \tag{2.3.3.1}
\]
et on a par suite, pour deux éléments $f$, $g$ de $S$
(\hyperref[I.1.1.9.1-fr]{I, 1.1.9.1})
\[
  \label{II.2.3.3.2-fr}
  D_+(fg)=D_+(f)\cap D_+(g).
  \tag{2.3.3.2}
\]
\end{env}

\begin{proposition}[2.3.4]
\label{II.2.3.4-fr}
Lorsque $f$ parcourt l'ensemble des éléments homogènes de $S_+$, les
$D_+(f)$ forment une base de la topologie de
$X=\operatorname{Proj}(S)$.
\end{proposition}

\begin{proof}
Il résulte en effet de (\hyperref[II.2.3.2.2-fr]{2.3.2.2}) et
(\hyperref[II.2.3.2.4-fr]{2.3.2.4}) que toute partie fermée de $X$ est
intersection d'ensembles de la forme $V_+(f)$, où $f$ est homogène de degré
$>0$.
\end{proof}

\begin{env}[2.3.5]
\label{II.2.3.5-fr}
Soit $f$ un élément \emph{homogène} de $S_+$, de degré $d>0$~; pour tout
idéal premier gradué $\mathfrak{p}$ de $S$, ne contenant pas $f$, on sait
que l'ensemble des $x/f^n$, où $x\in\mathfrak{p}$ et $n\geq 0$, est un
idéal premier de l'anneau de fractions $S_f$
(\hyperref[0.1.2.6-fr]{0, 1.2.6})~; sa trace sur $S_{(f)}$ est donc un idéal
premier de cet anneau, que nous désignerons par
$\psi_f(\mathfrak{p})$~: c'est l'ensemble des $x/f^n$ pour $n\geq 0$,
$x\in\mathfrak{p}\cap S_{nd}$. On a ainsi défini une application
\[
  \psi_f:D_+(f)\longrightarrow\operatorname{Spec}(S_{(f)})~;
\]
en outre, si $g\in S_e$ est un second élément homogène de $S_+$, on a un
diagramme commutatif
\[
  \label{II.2.3.5.1-fr}
  \xymatrix{
    D_+(f)\ar[r]^{\psi_f}
    & \operatorname{Spec}(S_{(f)})\\
    D_+(fg)\ar[u]\ar[r]_{\psi_{fg}}
    & \operatorname{Spec}(S_{(fg)})\ar[u]
  }
  \tag{2.3.5.1}
\]
où la flèche verticale de gauche est l'inclusion, et celle de droite est
l'application ${}^a\omega_{fg,f}$ déduite de l'homomorphisme canonique
$\omega=\omega_{fg,f}:S_{(f)}\to S_{(fg)}$
(\hyperref[I.1.2.1-fr]{I, 1.2.1}). En effet, si
$x/f^n\in\omega^{-1}(\psi_{fg}(\mathfrak{p}))$, où
$fg\notin\mathfrak{p}$, on a par définition
$g^n x/(fg)^n\in\psi_{fg}(\mathfrak{p})$, donc
$g^n x\in\mathfrak{p}$ et par suite $x\in\mathfrak{p}$, et la réciproque
est évidente.
\end{env}

\begin{proposition}[2.3.6]
\label{II.2.3.6-fr}
L'application $\psi_f$ est un homéomorphisme de $D_+(f)$ sur
$\operatorname{Spec}(S_{(f)})$.
\end{proposition}

\begin{proof}
En premier lieu, $\psi_f$ est continue~; car si $h\in S_{nd}$ est tel que
$h/f^n\in\psi_f(\mathfrak{p})$, on a par définition
$h\in\mathfrak{p}$ et réciproquement, donc
\[
  \psi_f^{-1}(D(h/f^n))=D_+(hf)
\]
et notre assertion résulte de la formule
(\hyperref[II.2.3.3.2-fr]{2.3.3.2}). En outre, les $D_+(hf)$, où $h$
parcourt les ensembles $S_{nd}$, forment une base de la topologie de
$D_+(f)$, en vertu de (\hyperref[II.2.3.4-fr]{2.3.4}) et de la formule
(\hyperref[II.2.3.3.2-fr]{2.3.3.2})~; ce
